% 广东工业大学本科生毕业设计(论文)LaTeX模板
% 根据《广东工业大学本科生毕业设计(论文)格式规范》制作

\documentclass[12pt,a4paper]{article}
\usepackage[UTF8]{ctex} % 中文支持
% 字体选择：仅在 XeLaTeX 或 LuaLaTeX 下使用 fontspec
\usepackage{ifxetex,ifluatex}
\ifnum\ifxetex1\else0\fi=1
    \usepackage{fontspec}
\else
    \ifnum\ifluatex1\else0\fi=1
        \usepackage{fontspec}
    \else
        \usepackage{times} % 对于 pdfLaTeX 保持传统字体设置
    \fi
\fi
\usepackage{geometry} % 页面设置
\usepackage{circuitikz} % 电路图支持
\usepackage{titlesec} % 标题格式设置
\usepackage{fancyhdr} % 页眉页脚设置
\usepackage{setspace} % 行距设置
\usepackage{graphicx} % 图片支持
\usepackage{amsmath} % 数学公式
\usepackage{amssymb} % 数学符号
\usepackage{booktabs} % 表格美化
\usepackage{array} % 表格支持
\usepackage{multirow} % 表格合并
\usepackage{longtable} % 长表格
\usepackage{cite} % 参考文献引用
\usepackage{enumitem} % 列表环境设置
\usepackage{indentfirst} % 首行缩进
\usepackage{float} % 浮动体控制
\usepackage{caption} % 图表标题
\usepackage{subcaption} % 子图
\usepackage{hyperref} % 超链接
\usepackage{tikz}
\usetikzlibrary{positioning, arrows.meta, calc, backgrounds, fit, shapes.misc, shapes.geometric} % TikZ 库
% 引用显示为方括号内的上标（整体上标），在 hyperref 之后定义为稳健命令
\makeatletter
\DeclareRobustCommand{\brsupcite}[1]{\textsuperscript{[\citen{#1}]}}
\makeatother
\let\oldcite\cite
\renewcommand{\cite}[1]{\brsupcite{#1}}
\usepackage{pdfpages} % 插入PDF封面
\usepackage{listings} % 代码块支持
\usepackage{xcolor} % 代码高亮支持
\usepackage{titlesec}
\titleclass{\subsubsubsection}{straight}[\subsubsection]

% 设置格式
\titleformat{\subsubsubsection}
  {\normalfont\normalsize\bfseries}{\thesubsubsubsection}{1em}{}
\titlespacing*{\subsubsubsection}{0pt}{3.25ex plus 1ex minus .2ex}{1.5ex plus .2ex}

% 设置编号
\newcounter{subsubsubsection}[subsubsection]
\renewcommand\thesubsubsubsection{\thesubsubsection.\arabic{subsubsubsection}}
% 代码块设置
\lstset{
    language=Verilog,
    basicstyle=\ttfamily\scriptsize,
    keywordstyle=\color{blue}\bfseries,
    stringstyle=\color{red},
    commentstyle=\color{green!60!black},
    numbers=left,
    numberstyle=\tiny\color{gray},
    frame=single,
    framesep=5pt,
    breaklines=true,
    breakatwhitespace=false,
    tabsize=2,
    showstringspaces=false,
    captionpos=b,
    abovecaptionskip=10pt,
    belowcaptionskip=10pt
}
\tikzstyle{startstop} = [rectangle, rounded corners, minimum width=3cm, minimum height=1cm,text centered, draw=black]
\tikzstyle{io} = [trapezium, trapezium left angle=70, trapezium right angle=110, minimum width=3cm, minimum height=1cm, text centered, draw=black]
\tikzstyle{process} = [rectangle, minimum width=3cm, minimum height=1cm, text centered, draw=black]
\tikzstyle{decision} = [diamond, minimum width=3cm, minimum height=1cm, text centered, draw=black]
\tikzstyle{arrow} = [thick,->,>=stealth]

% 页面设置
\geometry{
    top=30mm,
    bottom=25mm,
    left=30mm,
    right=20mm
}

% 公式编号格式：按章节顺序标注，如 4-1, 2-2
\numberwithin{equation}{section}
\renewcommand{\theequation}{\thesection-\arabic{equation}}

% 行距设置：1.5倍行距
\onehalfspacing

% 字体设置（为确保 \textbf 生效，显式设置粗体字体）
% 字体设置（若使用 XeLaTeX/LuaLaTeX 则设定 fontspec 字体）
\ifnum\ifxetex1\else0\fi=1
    \setmainfont{Times New Roman}[BoldFont={Times New Roman}, ItalicFont={Times New Roman Italic}]
    \setCJKmainfont[BoldFont={SimHei}]{SimSun} % 宋体，粗体使用 SimHei
    \setCJKfamilyfont{hei}{SimHei} % 黑体
\else
    % pdfLaTeX: 保持默认字体设置，\textbf 在常规字体下可用
\fi
\newcommand{\hei}{\CJKfamily{hei}}

% 字号定义
\newcommand{\erhao}{\zihao{2}}
\newcommand{\sanhao}{\zihao{3}}
\newcommand{\sihao}{\zihao{4}}
\newcommand{\xiaosi}{\zihao{-4}}
\newcommand{\xiaowu}{\zihao{-5}}

% 表格注释命令
\newcommand{\note}[1]{\par\vspace{0.5em}\noindent\small #1}

% 表格和图像按章节编号设置
\renewcommand{\thetable}{\thesection.\arabic{table}}
\renewcommand{\thefigure}{\thesection.\arabic{figure}}

% 每章开始时重置表格和图像计数器
\makeatletter
\@addtoreset{table}{section}
\@addtoreset{figure}{section}
\makeatother

% 表格标题格式设置：五号黑体加粗，表序与表名之间空一格
\captionsetup[table]{
  font={bf, small},
  labelfont={bf, small},
  labelsep=space,
  justification=centering
}

% 图像标题格式设置：五号宋体，图号与图名之间空一格
\captionsetup[figure]{
  font={small},
  labelfont={bf, small},
  labelsep=space,
  justification=centering
}

% 标题格式设置
\titleformat{\section}{\centering\hei\bfseries\sanhao}{\thesection}{1em}{}
\titleformat{\subsection}{\hei\bfseries\sihao}{\thesubsection}{1em}{}
\titleformat{\subsubsection}{\hei\bfseries\xiaosi}{\thesubsubsection}{1em}{}
\titleformat{\subsubsubsection}{\hei\bfseries\xiaosi}{\thesubsubsubsection}{1em}{}

% 增加标题深度支持
\setcounter{secnumdepth}{4}
\setcounter{tocdepth}{3}

% 段落格式
\setlength{\parindent}{2em} % 首行缩进2字符
\setlength{\parskip}{0pt} % 段间距为0

% 页眉页脚设置
\pagestyle{fancy}
\fancyhf{}
\fancyfoot[R]{\xiaowu\thepage} % 页码位于页脚右侧，小五号
\renewcommand{\headrulewidth}{0pt} % 去掉页眉横线

% 封面设置
% 封面设置（仅当封面文件存在时插入，避免编译因找不到文件失败）
\newcommand{\makecover}{%
      \IfFileExists{论文封面.pdf}{\includepdf[pages=-]{论文封面.pdf}}{}%
}

% 中英文摘要环境
\newenvironment{abstractcn}{
    \begin{center}
        \hei\bfseries\sanhao 摘要
    \end{center}
    \addcontentsline{toc}{section}{摘要}
}{\par}

\newenvironment{abstracten}{
    \begin{center}
        \textbf{\sanhao Abstract}
    \end{center}
    \addcontentsline{toc}{section}{Abstract}
}{\par}

% 关键词命令
\newcommand{\keywords}[1]{
    \vspace{\baselineskip}
    {\hei\bfseries\sihao 关键词：} #1
}

\newcommand{\keywordsen}[1]{
    \vspace{\baselineskip}
    {\textbf{\sihao Keywords:}} #1
}

% 目录格式设置
\usepackage[titles]{tocloft}
\renewcommand{\cftsecfont}{\hei\bfseries\xiaosi} % 一级标题小四黑体
\renewcommand{\cftsubsecfont}{\songti\xiaosi} % 二级标题小四宋体
\renewcommand{\cftsubsubsecfont}{\songti\xiaosi} % 三级标题小四宋体
% 文献翻译元信息
\title{外文文献翻译}
\author{请填写姓名}
\date{请填写日期}
\providecommand{\tightlist}{%
  \setlength{\itemsep}{0pt}\setlength{\parskip}{0pt}}
\providecommand{\passthrough}[1]{#1}
\definecolor{shadecolor}{RGB}{248,248,248}
\newenvironment{Shaded}{}{}


\begin{document}

% 文献翻译封面：仅插入第 1 页，避免多页封面 PDF 的异常页导致编译失败。
\IfFileExists{g:/My_code/Agent_atuo_graduation_project_of_GDUT/模板/文献翻译封面.pdf}{\includepdf[pages=1]{g:/My_code/Agent_atuo_graduation_project_of_GDUT/模板/文献翻译封面.pdf}}{}

\pagenumbering{gobble}
\tableofcontents
\clearpage
\pagenumbering{arabic}
\setcounter{page}{1}


\section{外文文献译文（1）}
\subsection{译文信息}
\begin{itemize}
\item 外文题名：Analysis and Comparison of UART, SPI and I2C
\item 中文译名：UART、SPI 与 I2C 的分析与比较
\item 作者：未在原文中明确给出
\item 来源/期刊/会议/出版社：未在原文中明确给出
\item 年份：未在原文中明确给出
\item DOI 或 URL：未在原文中明确给出
\end{itemize}

\subsection{翻译正文}
\subsubsection*{摘要}

当前可选的硬件通信协议较多，其中 UART、SPI 和 I2C 是最具代表性的三类协议，并且已经得到广泛应用。它们各自具有不同特点，因此如何根据实际需求选择合适协议，长期以来都是设计人员需要面对的问题。本文分别介绍三种协议的组成方式、工作原理以及优缺点。通过比较三者的传输距离可以发现，目前 SPI 协议在特定条件下能够支持更长距离传输；在传输速率方面，SPI 协议具有更高速度，但受到较多限制条件约束；UART 的性能较为稳定；SPI 还提供多种传输模式可供选择。在传输功耗方面，UART 具有明显优势，优于另外两种协议。本文希望从设计目标出发，为协议选择提供一种新的思考方式，并由此形成观察协议选型的新视角。

关键词：UART，SPI，I2C

\subsubsection*{引言}

目前存在多种通信协议，每种协议都有自身的优点和不足，在不同应用场景下，硬件传输协议的选择也会有所差异。本文讨论的主题是如何选择合适的协议。下文将讨论当前使用较广泛的几种通信协议，并分析多种协议的特性，从而判断哪些协议在特定场景下更合适、更具优势。UART 是一种常用的硬件通信协议，只需要两条数据线即可完成双向数据传输，并且不需要时钟信号进行同步操作。它以数据包形式传输数据，每个数据包的数据长度可以变化。为了保证数据传输过程中的安全性和完整性，数据包中可以加入校验位。发送端和接收端通过比特率实现同步，因此两个工作频率不同的设备也能够通信。

SPI 协议是一种速度较快的传输协议，允许主设备以较高频率与从设备通信。SPI 协议只允许存在一个主设备，但可以连接多个从设备。主设备负责产生用于与从设备同步的时钟信号。SPI 在芯片中通常只占用四根线，能够节省引脚资源，也有利于节省 PCB 布局空间。I2C 协议是一种两线串行总线，只需要串行数据线和串行时钟线两根双向信号线。总线接口可以集成在芯片内部，从而优化空间和成本。它支持多主设备和多从设备，但同一时刻只允许一个主设备控制总线。连接到 I2C 总线上的设备数量主要受总线最大电容限制。I2C 具有低功耗和较强抗干扰能力等特点。通过理解上述协议的工作原理与特点，可以在特定工作环境或一般工作环境中得到更合适的方案，例如面向低功耗场景、长距离传输需求或高速传输需求进行协议选择。本文旨在为协议搭配和选择提供新的建议。

\subsubsection*{UART}

通用异步收发传输器（Universal Asynchronous Receiver/Transmitter，UART）是一种异步串行通信协议，也是最常用的硬件间通信协议之一。当前，UART 已应用于多种场景，例如各类环境传感器、无线通信模块等。

\paragraph{UART 的工作原理}

UART 是一种串行通信协议，不同于并行通信，其传输数据不会在同一时刻全部到达接收端，而是通过一条数据线逐位发送给接收端。这意味着它的传输速率可能低于并行传输协议。发送端的 TX 引脚直接连接接收端的 RX 引脚。数据在 UART 与控制设备之间以并行方式传递，控制设备可以是微控制器或微处理器。当通过 TX 引脚发送数据时，发送端 UART 需要将并行信息转换为串行数据并发送给接收端。接收端接收该数据后再转换回并行形式，并与其控制器通信。

UART 区别于其他通信协议的另一个特点是，发送端和接收端并不通过时钟信号同步，而是以相同波特率实现同步。波特率用每秒传输的比特数表示，单位为 bps，常见取值为 9600 或 115200。

\paragraph{UART 的结构}

UART 传输的数据以数据包形式组织。每个数据包包含 1 个起始位、5 至 9 个数据位、1 个奇偶校验位（可选择是否使用），以及 1 个、2 个或 1/2 个停止位。

起始位：当数据线空闲时，UART 协议处于高电平状态。当数据传输开始时，发送端会在一个时钟周期内将传输线电平由高变低。起始位被添加在实际数据之前。接收端检测到数据线从高到低的变化后，开始接收实际数据。

数据位：数据位是发送端与接收端之间实际传输的数据。数据长度可以为 5 至 9 位。特别地，不使用奇偶校验位时可以得到 9 位数据长度。多数情况下，最大数据长度为 8 位。

奇偶校验位：奇偶校验位可以帮助接收端检查传输过程中数据是否发生改变。数据位可能因为多种错误而改变，例如接收溢出导致数据丢失，或者接收端接收了大量数据但无法一次处理。未完成的 UART 传输会造成数据丢失，也可能与发送后进入休眠、接收设备断电等情况有关。接收端读取数据后，会统计数据中 1 的个数，并判断总数为奇数还是偶数。如果校验位为 0，表示数据包中 1 的个数为偶数；如果校验位为 1，表示数据包中 1 的个数为奇数。

停止位：顾名思义，停止位用于标记数据包的结束。通常停止位占用 1 个或 2 个时钟周期，也可以设置为 1/2 个时钟周期。当停止位结束后，数据线进入空闲状态，其电平回到高电平。

\paragraph{UART 的特点}

与其他硬件通信协议不同，UART 只需要两条数据线即可完成全双工数据传输，而 SPI 协议至少需要四条数据线互连。UART 也不需要时钟线或其他时间信号来完成两个设备之间的通信。由于存在奇偶校验位，基本的数据传输准确性可以得到保证。UART 具有广泛的硬件应用和多种使用场景。该协议还具有良好的功耗表现，有助于绿色通信技术的发展，并已在 Artix-7 等平台中得到验证。

UART 的不足主要包括以下几点：单个数据包的最大传输长度限制为 9 位；与并行传输相比，UART 的传输速率较低；UART 协议不是总线通信协议，不能同时与多个设备通信，同一时刻只能进行一对一通信。与 SPI 协议中一个主机可连接多个从机、I2C 协议中可存在多个主机和多个从机不同，UART 通信的两个设备之间波特率差异通常不应超过 10%。

\subsubsection*{SPI}

SPI 是串行外设接口（Serial Peripheral Interface）的缩写。它是一种全双工、高速同步通信总线，通常只使用芯片上的四个引脚，因此能够节省芯片引脚、释放空间并简化 PCB 布局。由于该协议使用较为简单，越来越多的芯片开始集成这种通信协议，例如 AT91RM9200。其工作原理如原文图 1 所示。

SPI 工作在主从模式下。从设备通常可以是 EPROM、Flash、AD/DA、音视频处理芯片或其他器件，而主机通常是 FPGA、MCU 或 DSP 等可编程控制器。SPI 主要由 SCLK、CS、MOSI 和 MISO 等信号线组成。SCK、SS、SDI 和 SDO 等名称也具有相同含义。当存在多个从设备时，CS 用于选择需要控制的从设备。

\paragraph{SPI 的工作原理}

一般情况下，器件的工作模式是固定的，主机需要采用相同工作模式。控制器通常指 SPI 设备中的控制寄存器，可以通过配置控制寄存器设置 SPI 总线的传输模式，使双方能够正常通信。

当 SPI 空闲时，时钟信号的极性 CPOL 表示时钟处于高电平还是低电平。当器件未使用且 CPOL 设为 1 时，SCK 引脚上的时钟信号为高电平；当 CPOL 设为 0 时，情况相反，即 SCK 空闲时为低电平。CPOL=0 表示 SCK 未使用时为 0，CPOL=1 表示 SCK 未使用时为 1。

CPHA 表示时钟相位，用于指定 SPI 设备是在 SCK 引脚时钟信号跳变到上升沿还是下降沿时开始采样数据。如果 CPHA 设为 1，时钟信号的上升沿和下降沿会分别使 SPI 设备采样数据和传输数据。当 CPHA 设为 0 时，情况相反。CPHA=0 表示输入和输出数据在 SCK 的第一个边沿有效；CPHA=1 表示输入和输出数据在 SCK 的第二个边沿有效。在四种模式中，模式 0 和模式 3 使用最广，大多数 SPI 设备都支持这两种模式。

\paragraph{SPI 的结构}

默认情况下，讨论 SPI 时通常指传统的四线 Motorola SPI 协议，它总共有四条数据线：SCLK、MOSI、MISO 和 CS。典型四线系统的优点是能够实现全双工数据传输。一个 SPI 接口只能有一个主机，但可以有一个或多个从机。只有一个主机和一个从设备时需要一条 CS 线，而多个从设备则需要多条 CS 线。MOSI 和 MISO 是数据线，其中 MOSI 用于主机向从机发送数据，MISO 用于从机向主机发送数据。SCLK 是时钟信号。主机输出、从机输入、主机数据 MISO、主机输入、从机输出、CS 从设备选择以及低电平有效等信号共同构成 SPI 通信，SPI 速率由时钟频率决定。连接示意图如原文图 2 所示。

根据不同应用场景，三线 SPI 主要有两种类型。对于部分通信外设，其只有三条 SPI 总线，即一条 CS、一条 CLK 和一条 MOSI/MISO，这意味着数据输入与输出共用同一条数据线。此时仅存在 SCLK、SDIO 和 CS 三根线，SDIO 是用于半双工通信的双向端口。例如，ADI 的许多 ADC 芯片支持双向传输。双向端口在作为输入使用时，应由 FPGA 设置为高阻态 Z。

\paragraph{补充说明}

SPI 总线具有若干特殊特性，不仅限于其高频通信速度，还包括菊花链连接方式，这些特性说明 SPI 是一种具有优势的通信协议。

第一，不同的连接结构。SPI 有两种连接方式。第一种是传统的独立从机结构，每个从机都需要自己的 CS 线。当主机希望与某个从机通信时，会拉低对应 CS 信号线，同时保持其他 CS 信号线为高电平。由于未被选中的从机的 MISO 引脚与其他从机共用同一信号线，因此这些从机的 MISO 引脚必须配置为高阻输出。

第二种连接方式是菊花链配置。菊花链是指信号线从一个设备串行传输到下一个设备，直到数据到达目标设备。

菊花链的主要缺点是，如果某个从机发生单点故障，则优先级低于该设备的从机可能会断开。距离主机更远的从机，其服务优先级也会更低。如有必要，应对后续优先级和从机线路检测器进行配置。

为了避免单点故障导致整个链路失效，应在从机服务器超时时快速响应。

在菊花链模式下，每个从机会在下一个时钟周期将输入数据复制到输出端，从而充分利用 SPI 移位寄存器。差异比较图如原文图 3 所示。

第二，传输速率。与 I2C 的标准模式 100K、快速模式 400K 和高速模式 3.4M 不同，SPI 协议没有固定速率。SPI 的最大时钟频率、CPU 解析 SPI 数据的能力以及输出端驱动能力共同决定 SPI 传输速率，且最大信号传输速率还应遵循 PCB 布局要求。

\subsubsection*{I2C}

I2C 总线也可称为 IIC 或 Inter-Integrated Circuit。该名称中的数字 2 更接近平方式写法，而不是普通数字 2。I2C 是 Philips Semiconductor 于 1982 年提出的一种同步、多控制器、多目标串行通信总线。这种串行通信总线最重要的特点是，可以用两根线连接多个设备，且设备有机会在主设备和从设备角色之间转换。此外，开发者还可以选择不同通信模式。I2C 用于在嵌入式系统中连接微控制器、EEPROM、I/O 接口以及其他外设。例如，Arduino 与气压传感器之间的连接就是 I2C 总线的典型应用。ADC 和 DAC 芯片可以将模拟数据转换为数字 I2C 信号，从而保持数值稳定和准确。

\paragraph{I2C 的工作原理}

I2C 使用 SCL 和 SDA 两根线连接多个设备，它们分别称为串行时钟线和串行数据线。每个设备拥有一个 7 位地址码。根据 i2c-bus.org 的说明，7 位宽度理论上允许 128 个 I2C 地址，但由于某些原因，其中一些地址具有特殊用途，例如 0000000 0、0000000 1、0000001 X、0000010 X、0000011 X、00001XX X、11110XX X 和 11111XX X，这些均以二进制形式表示，并与 10 位地址相关。

I2C 仿真和概念基于方波振荡，其中没有考虑电磁干扰或其他干扰因素。通过在 SCL 和 SDA 线上使用电容，可以降低和滤除波形噪声。然而，电容特性和充电周期会受到电容大小影响。因此，选择合适电容值非常重要，因为时钟线和数据线必须达到逻辑高电压，否则信息将无法被识别。根据 Texas Instruments 的建议，I2C 总线上每个设备使用的电容在标准模式和快速模式下应遵循 400 pF 的限制，在 Fast Mode Plus 下规定为 550 pF。相关波形如原文图 4 所示。

\paragraph{I2C 的结构}

当 I2C 总线空闲时，SCL 和 SDA 两条线均保持高电平。当设备准备传输数据时，SDA 会切换为低电平并开始通过 SDA 线传输数据，时钟线则向各个设备提供时钟信息。当设备准备停止传输数据时，SDA 会在 SCL 为高电平时由低电平切换为高电平。

在写数据时，设备首先向总线发送停止信号，以降低写入失败的情况；随后发送复位信号和起始信号。接着发送 7 位地址，从设备会返回反馈信号。之后开始传输 8 位数据，并在每个周期从从设备获得反馈信号。结束时，向 I2C 总线发送停止信号，并释放总线进入空闲状态。

\paragraph{I2C 的不足}

由于多种原因，I2C 不适合用于长距离通信。虽然存在一些方法可以解决位翻转问题，但来自 PCB 或导线的寄生电容属于物理缺陷，难以完全消除。如果导线材料较差，寄生电容和导线电阻的波动会进一步影响 I2C 总线稳定性。如果电阻过大，主设备将无法解析从设备数据，从而引发一系列可靠性问题。

\subsubsection*{比较}

前面章节介绍了三种协议的工作原理和特点，本章从单一方向对它们进行比较。在传输距离方面，三种协议各自都有长度限制。UART 协议的一般传输距离为 15 m；但在基于 RS-485 的设计中，可以实现 100 m 的长距离传输。SPI 协议通常属于短距离传输协议，但在特殊设计环境下可实现超过 1200 m 的传输距离。I2C 是短距离传输协议，在不使用中继器的情况下，其传输距离约为 200 mm 至 300 mm。从上述传输特性可以看出，在应用环境需要长距离传输时，可以考虑使用 SPI 协议。

UART 协议在不同标准和不同传输速率下工作，其在 RS-232 标准下的传输速率为 20 Kbps。在 RS-422 标准下，传输速率可以达到 10 Mbps。同时，传输距离越长，最大传输速率越小。I2C 协议在不同模式下具有不同最大传输速率，其在标准模式下最高为 100 Kbps，在超快速模式下最高可达 5 Mbps。SPI 协议的传输速率最高可达 50 Mbps，但其最大传输速率受三个条件限制：系统时钟频率、CPU 处理 SPI 数据的能力以及 PCB 允许的最大信号传输速率。如果设计需要高速传输，SPI 协议的高传输速率是一种选择；如果设计要求连接简单和布局简洁，I2C 可以满足需求；如果设计目标是节能，则 UART 是最佳选择，其功耗为 0.0135 W，明显低于另外两种协议。

\subsubsection*{结论}

在本研究中，通过学习 UART 的工作原理、格式描述和特性可以发现，UART 的传输速率可能低于某些并行协议，发送端和接收端不依赖其他介质进行同步，并以数据包形式完成传输。最重要的是，UART 数据传输只需要两条数据线，并且是全双工系统。虽然它只能进行一对一通信，但也可以通过多个从机实现一对多的应用形式。

随后，本文从多个方面分析 SPI 协议，并得出结论：SPI 作为一种全双工高速通信线路，可以简化布局，主机也可以自行改变控制器。四线结构是全双工，而三线结构只是半双工。它可以配置主传输模式，速率也会不断变化。菊花链连接无疑是一种较好的应用方式，但其缺点是容易导致单点故障。

最后，本文回顾了 I2C 的相关文献并形成了认识。I2C 可以用两根线连接多个设备，其特性可以发生变化，并且能够在一定程度上降低波形噪声。与此同时，I2C 对导线材料要求较高，因为导线材料会影响总线稳定性，过大的电阻也会影响服务器数据解析。


\clearpage
\section{外文文献原文（1）}
\subsection{文献原文}
\subsubsection*{Analysis and Comparison of UART, SPI and I2C}

\paragraph{Abstract}

Nowadays there are many hardware communications protocols to choose
from, UART, SPI, I2C three protocols are the most representative, they
are widely used. They have different characteristics, and how to choose
the right choice is a problem that has been bothering people for a long
time. In this article, the composition of the three protocols, how they
work, and their advantages or disadvantages are described separately.
And through the comparison of their transmission distance, it will be
found that today\textbackslash{}'s SPI protocol can transmit a longer distance, in the
comparison of transmission rates SPI protocol has a faster rate but many
limited conditions, UART has a more stable performance, SPI has a
variety of transmission modes to choose from. In terms of transmission
power consumption, UART has a huge advantage, outperforming the other
two protocols. Through the article, we hope to provide a new way of
thinking from the purpose required for the design. This leads to a new
perspective on protocol selection

Keywords：UART, SPI, I2C

\paragraph{INTRODUCTION}

Various communication protocols are available nowadays, each protocol
has its own advantages and disadvantages, and the choice of hardware
transmission protocol is also different in different situations. The
topic to be covered in this post is how to pick a suitable protocol. The
following article discusses a few of the most well-liked communication
protocols now in use and examines numerous protocols\textbackslash{}' properties to
determine which are preferable and superior. UART is a widely used
hardware communication protocol that requires only two data lines to
complete bidirectional data transmission and does not require

clock signals for synchronous operation. It transfers data in packets,
each with a variable size of data. A check digit can be included in the
package to guarantee security and integrity during data transmission.
Synchronization between the transmitter and receiver is achieved by bit
rate, which allows two devices operating at different frequencies to
communicate. The SPI protocol is a very fast transmission protocol that
allows the master to communicate with the slave at a higher frequency.
The SPI protocol allows only one master to exist, but multiple slaves
can exist. The master is responsible for generating a clock signal that
is used to synchronize with its slave. The SPI saves pins by occupying
only 4 wires in the chip, which also benefits saving PCB layout space.
The I2C protocol is a two-wire serial bus that requires only two
bidirectional lines, a serial data line, and a serial clock line. The
bus interface is formed within the chip, optimizing space and cost. It
supports multi-master and multi-slave devices, but only one host is
allowed at a time. The number of I2C protocols connected to the bus is
limited only by the maximum capacitance of the bus. It is characterized
by low power consumption and strong antiinterference ability. Through
the working principle and characteristics of the above protocols, we get
a better solution in the specific working environment or the general
working environment and provide new solutions for the application of the
protocol in different scenarios, in the pursuit of lowpower environment
solutions or the choice of long-distance transmission, or the pursuit of
excellent choices under highspeed transmission. The entire article aims
to offer fresh suggestions for protocol collocation and selection.

\paragraph{UART}

Universal Asynchronous Receiver/Transmitter is an Asynchronous serial
communication protocol. It is one of the most widely used
hardware-to-hardware communication protocols. Nowadays UART has been
used in many applications such as various types of environmental
sensors, wireless communication modules, and so on.

A.  Working principle of UART

UART is a serial communication protocol, that differsfrom parallel
communication, transmissions in that the transmitted data does not reach
the receiver at the same time. It is transmitted bit by bit to the
receiver via a data line. This means that its transfer rate may be
slower than the transfer rate of the parallel transfer protocol. The
transmitter's TX pin is directly connected with the receiver's RX pin.
Data travels to and from a UART in parallel to the control device, such
as microcontrollers and microprocessors. When sending on the TX pin, the
transmitter's UART needsto translate parallel information into serial
and transmit it to the receiver. The receiver receives this data and
translates it back to parallel, then communicate with its controller.
Another feature that distinguishes UART from other communication
protocols is that the transmitter and receiver are not synchronized
through the clock signal. Instead, the transmitter and receiver
synchronize with each other at the same baud rate. The baud rate is
expressed in bits per second or bps, it is usually 9600 or 115200.

B.  Structure of the UART

    UART transmitted data is organized into packets. Each packet
    contains 1 start bit, 5 to 9 data bits, a parity bit (ability to
    choose whether to use or not), 1 or 2 or 1/2 stop bits.

    START BIT: The UART protocol is in a high-level state when the data
    line is idle. When data transfer starts, the transmitter will change
    the level of the transmission line from high to low for one clock.
    It is added before the actual data. The receiver detects this change
    from high to low on the data line and starts receiving the actual
    data.

    DATA BITS: The data bits are the actual data that is transmitted
    between transmitter and receiver. The length of data can be anywhere
    between 5 to 9 bits. In particular, 9-bit data is the resulting
    length without the use of parity bits. Most of the time the maximum
    data length is 8 bits.

    PARITY BIT: The parity bit can help the receiver check if any data
    has changed during the transmission. Bits can be changed in kinds of
    error, such as receiving overflow lost data, the reason it has
    received lots of data but it cannot process at once. Unfinished UART
    transmission results in data loss, use this function to enter
    hibernation after sending, power off the receiving device, and so
    on. After the receiver reads the data, the number of digits is
    counted as 1 and determines whether the total is odd or even if the
    parity bit is 0, the number of 1s in the packet is even, and if the
    parity bit is 1, the number of 1s in the packet is odd.

    END BITS: As the name suggests, it is marking the end of the data
    packet. Typically, the stop bit occupies 1 or 2 clocks, but it is
    also possible to set 1/2 clocks. When the stop bit ends, the data
    line is idle and its level returns to the high position.

C.  Features of UART

    Unlike other hardware communication protocols, UART only requires
    two data lines for full-duplex data transmission, such as the SPI
    protocol, which requires at least four data lines to be connected to
    each other. UART also does not require a clock or other time signal
    for communication between two devices. Because of the presence of a
    parity bit, the basic accuracy of the transmitted data is ensured
    when used. UART has a wide range of hardware applications and a
    variety of use scenarios. The UART protocol also has excellent power
    performance, enabling an advance in the development of green
    communication technology and is proven in Artix-7.

    As for its disadvantages, there are mainly the following points: The
    maximum package size transmitted is limited to 9 bits. Compared to
    parallel transmissions, UART has a lower transmission rate. The UART
    protocol is not a bus communication, it cannot communicate with
    multiple devices at the same time, at a time can only be one-to-one
    communication. Unlike the SPI protocol, a host can have multiple
    slaves, the I2C protocol can allow multiple hosts to exist with
    multiple slaves. The difference in baud rate between two UART
    protocol devices that communicate with each other should not exceed
    10\%.

\paragraph{SPI}

The serial Peripheral Interface is known as SPI. It is a full-duplex,
high-speed synchronous communication bus that only uses four pins on the
chip, saving the chip\textbackslash{}'s pins, freeing up space, and making the PCB
layout easier. More and more chips, including AT91RM9200, are
integrating this communication protocol because of how simple it is to
use.The working principle is shown in Figure 1 below.

\noindent\textit{[原文图片位置：图片需根据 DOCX 原文人工补齐。]}
height="1.7104166666666667in"\}

In master-slave mode, SPI operates. The slave is typically an EPROM,
Flash, AD/DA, audio and video processing chip, or other devices, while
the host is typically a programmable controller such as an FPGA, MCU, or
DSP. The SCLK, CS, MOSI, and MISO lines make up the majority of them.
SCK, SS, SDI, and SDO are examples of names that all have the same
meaning. When there are multiple slave devices, the CS is used to select
the slave device to be controlled.

A.  Working principle of SPI

    In general, the working mode of the machine parts is fixed, and the
    host machine needs to adopt the same working mode. The Controller
    generally refers to the control registers in the SPI device, which
    can be configured to set the transmission mode of the SPI bus, so
    that both sides can normally \textbackslash{}"communicate\textbackslash{}".

    When the SPI is idle, the clock signal\textbackslash{}'s polarity---CPOL---shows
    whether it is high or low. When the device is not in use, the clock
    signal at the SCK tube\textbackslash{}'s foot is high if CPOL is set to 1. When
    CPOL is set to 0, the opposite is accurate. SCK is zero when it is
    not in use, as shown by CPOL=0. When SCK is in use, CPOL=1 indicates
    that it is 1.

    CPHA: clock phase, which specifies whether the SPI device starts
    sampling data when the clock signal on the SCK pin transitions to a
    rising edge or a falling edge. If CPHA is set to 1, the clock
    signal\textbackslash{}'s rising and falling edges cause the SPI device to sample
    data and transfer data, respectively. When CPHA is set to zero, the
    opposite is accurate. At the first edge of SCK, CPHA=0 denotes that
    the input and output data are legitimate. The input and output data
    are valid at the second edge of SCK when CPHA=1. Of the four modes,
    modes 0 and 3 are the most widely used, and most SPI devices support
    both modes.

B.  Structure of the SPI

    By default, when we discuss SPI, we are referring to the
    conventional 4-wire Motorola SPI protocol, which has four data lines
    total: SCLK, MOSI, MISO, and CS. The typical 4-wire system has the
    benefit of enabling full duplex data transfer. An SPI interface can
    have only one host, but it can have one or more slaves. One CS is
    required when there is just one host and one slave device, however,
    many slave devices require several CS. MOSI and MISO are data lines.
    MOSI sends data from the master to the slave, and MISO sends data
    from the slave to the master. The beginning of every data line:
    clock signal, or SCLK. The host output, slave input, host data MISO,
    host input, slave output, CS, slave device selection, and active low
    level are all components of the SPI rate, which is determined by the
    clock frequency. The connection schematic was shown in figure 2
    below.

    ![6c9ec775-25e5-4335-b5af-eded3fa6d0f2](media/image2.png)\{width="3.4118055555555555in"
    height="1.7826388888888889in"\}

    There are two primary varieties of 3-wire SPI, depending on the
    various application scenarios:

    For simcom communication, some peripherals have only three SPI
    buses, one CS, one CLK, and one MOSI/MISO, which means that the
    input and output of data are the same data line. The SCLK, SDIO, and
    CS wires are the only ones present. A bidirectional port for
    half-duplex communication is the SDIO. For instance, ADI\textbackslash{}'s numerous
    ADC chips enable bidirectional transmission. Bidirectional ports
    should be operated by FPGAs with high resistance state Z set when
    utilized as input.

C.  Additional information

    There are several special characteristics shown throughthe SPI bus,
    not only limit to its high-frequencycommunication speed but also the
    daisy chain connection type can be shown that SPI is a superior
    communication protocol.

<!-- -->

1)  Different connection structure

    SPI has two types of connections. The traditional independent slave
    setup is the first. Each slave requires their own CS line. The
    matching CS signal line is brought down while the other CS signal
    lines are kept high when the host wants to talk to a certain slave.
    The MISO pins of the unselected slave machine must be set up with
    high-resistance output since they are on the same signal line as the
    slave machine\textbackslash{}'s MISO pins.

    Daisy chain configuration is SPI\textbackslash{}'s second connecting option. Daisy
    chaining is the process of serially transmitting signal wires from
    one device to the next until the data reaches the destination
    device.

    The primary drawback of the Daisy chain is that it will disconnect
    the slave machine that has a lower priority than the device if the
    slave machine has a single point of failure. The priority for
    service will be lower for the slave machine that is further away
    from the host machine. If necessary, configure the priority going
    forward and the slave line detector.

    To avoid a single point of failure causing the entire link to fail,
    respond quickly if a slave server times out.

    The SPI shift register is fully utilized in daisy-chain mode when
    each slave copies input data to output in the subsequent clock
    cycle. The difference comparison diagram is shown in Figure 3.

    ![cb003dbb-7b59-4579-9fcc-0865b881b7d0](media/image3.png)\{width="4.177083333333333in"
    height="1.6847222222222222in"\}

2)  Transmission rate

    In contrast to I2C\textbackslash{}'s normal mode of 100K, fast mode of 400K, and
    high-speed mode of 3.4m, the SPI protocol does not have a set rate.
    Instead, the maximum SPI clock frequency, the CPU\textbackslash{}'s capacity to
    interpret SPI data, and the driving capacity of the output end all
    play major roles in determining the SPI transmission rate (maximum
    signal transmission rate should be followed by the PCB layout).

\paragraph{I2C}

The full name of the I2C bus can be called IIC or InterIntegrated
Circuit. The number two in the phrase would likely be square, not two.
The I2C is a synchronous, multicontroller and targets serial
communication bus, developed by Philips Semiconductor in 1982. The most
significant point of this serial communication bus is that it can
connect with multiple devices with two lines, and the device can have
the chance to change its characteristic with master and slave. Also,
there are different communication modes that can be chosen by
developers. I2C is used to connect devices like microcontrollers,
EEPROMs, I/O interfaces, and other peripheral devices in an embedded
system. For example, air pressure sensor connected between Arduino and
the sensor would be a perfect application for the I2C bus. The ADC and
DAC chip would convert the analog data to a digital i2c signal to keep
the number constant and accurate.

A.  Working principle of I2C

    I2C uses two wires, SCL and SDA, also known as Serial Clock Line and
    Serial Data Line, to connect with dozens of devices. Each device
    owns an addressing code with 7 bits. From i2c-bus.org, it suggested
    that the bite width with 7 can theoretically allow 128 I2C
    addresses, but for certain reasons, there are several addresses with
    special purposes, which are "0000000 0, 0000000 1, 0000001 X,
    0000010 X, 0000011 X, 00001XX X, 11110XX X, 11111XX X", has been
    noted as binary, 10-bit addresses.

    The i2c simulation and concepts are based on the square wave
    oscillation, in which electromagnetic or other interference has not
    been considered. By using capacitors on both SCL and SDA wire could
    reduce and filter the noise of the wave. However, the character of
    the capacitor and the charging period would affect by the size of
    the capacitor. It is extremely important to consider what size of
    the capacitor to use since the clock and data line should reach the
    logical high voltage, or the massage wouldn't be recognized. From
    Texas Instruments' suggestion, the capacitor used for the I2C bus of
    each device should follow the 400 pF for both Standard Mode and Fast
    Mode and specifies 550 pF for Fast Mode Plus. The following figure 4
    shows the waveform.

    ![b9f78690-17fe-4845-b61d-3d02ed418221](media/image4.png)\{width="3.8493055555555555in"
    height="1.1506944444444445in"\}

B.  Structure of the I2C

    As the I2C bus is resting, both SCL and SDA wires arekeeping HIGH.
    When the device is planning to transfer data, SDA would switch to
    LOW and start transferring data through the SDA wire, and the clock
    wire would provide the clock message to each of the devices. When
    the device is planning to stop transferring data, the SDA would
    switch from LOW to HIGH while the SCL is in HIGH condition.

    For writing data, the device first sends the stop signal to the bus,
    reduces the condition of failing to write; then sends the rest
    signal, the start signal. Send the address of 7 bits, and the slave
    device would send back a feedback signal. Then start transferring
    data with 8 bits, getting back a feedback signal from the slave as a
    period. For ending, send a stop signal to the I2C bus, and release
    the bus in an idle state.

C.  Flaws of I2C

    For several reasons, there are some disadvantages thatI2C is not
    suggested to be applicated in long-distance communication. Though
    there are some ways that can solve the bit flip issues, the
    parasitic capacitance from the PCB or wire cannot be solved as a
    physical defect. If the material of the wire is inferior, the
    fluctuation of the parasitic capacitance and the resistance of the
    wire will further affect the stability of the i2c bus. If the
    resistance is too large, the master cannot parse the data of the
    slaves, which will lead to a series of unreliable problems.

\paragraph{COMPARISON}

The previous chapters describe the working principles and features of
the three protocols, and in this chapter they will be compared in a
single direction. In terms of transmission distance, each of the three
protocols has its own limited length, and the transmission distance of
the UART protocol is 15m in general. However, in the design based on
RS-485, it is possible to achieve a long-distance transmission of 100m.
The SPI protocol is typically a short-range transmission protocol that
can achieve transmission distances of more than 1200m in a specially
designed environment. I2C is a short-range transmission protocol whose
transmission distance is about 200mm to 300mm without the use of
repeaters. As can be seen in the above transmission characteristics, the
SPI protocol can be used under the requirements of the application
environment that requires long-distance transmission. The UART protocol
operates at different rates under different standards at different
transfer rates, with a transfer rate of 20Kbps under RS-232 standards.
Under the RS-422 standard, its transmission rate can reach 10Mbps. And
the longer the transmission distance, the smaller the
maximumtransmission rate. The transfer rate of the I2C protocol has
different maximum transfer rates in different modes, and its transfer
rate is up to 100Kbps in standard mode. In ultra-fast mode it has a
maximum transfer rate of 5Mbps. The transmission rate of the SPI
protocol can reach up to 50Mbps, but its maximum transmission rate is
subject to three conditions: 1. System clock frequency 2. CPU processing
SPI data capability 3. Maximum signal transmission rate allowed by the
PCB. If the design requires high-speed transmission, the SPI protocol\textbackslash{}'s
high rate is an option. If the design requires simple connections and a
streamlined layout, I2C can do it. If the design is designed to save
energy, then UART is the best choice, it consumes much less power than
the other two protocols at 0.0135W.

\paragraph{CONCLUSION}

In this experiment, after studying the working principle, format
description and personality characteristics of UART, we found that: The
transmission rate of UART may be lower than some parallel protocols, and
the transmitter and receiver are synchronized with each other without
any medium. And use packet form in transmission. The most important
thing is that UART data transmission only requires two data lines, and
it is a full duplex system, although it can only be one-toone, it can
have multiple slave machines to achieve one-to many.

Then, we analyse all aspects of the SPI protocol and conclude that SPI
as a full-duplex and high-speed communication line can simplify the
layout and the host can change the controller by itself. The fourth line
is a full duplex, while the third line is only half a duplex. It can be
configured to set the main transmission mode, and the rate is constantly
changing. Chrysanthemum link is undoubtedly a very good way of
application, but the disadvantage is easy to cause the failure of the
point.

Finally, we reviewed the relevant literature on I2C and made a
self-perception. I2C can connect multiple Settings with two lines, its
characteristics can be changed, and the wave noise can be reduced to a
certain extent. Finally, it has high requirements on wire material,
because it will affect the stability of the bus, and too much resistance
will affect the analysis of server data.

\end{document}
